%=====================================================================
\chapter{Divide and Conquer: Merge Sort and the Pattern}\label{ch:mergesort}
%=====================================================================
\locator{2}{Part II --- Recursion and Divide-and-Conquer}

\begin{outcomes}
By the end of this chapter you will be able to:
\begin{tight}
  \item \textbf{State} the three steps of the divide-and-conquer pattern and
        \textbf{apply} it to a stated problem \emph{(CLO-6)}.
  \item \textbf{Prove} merge sort correct from the merge step's invariant, and
        \textbf{derive} its $\bigTh{n\lg n}$ bound in all three cases.
  \item \textbf{Explain} why merge sort's bound is case-independent and
        \textbf{contrast} it with insertion sort's.
  \item \textbf{Measure} the effect of an insertion-sort cutoff and
        \textbf{choose} one for a given machine.
  \item \textbf{Account} for a measured gap between two runs of the same
        algorithm in terms of work done and hardware behaviour.
\end{tight}
\end{outcomes}

\begin{prereq}
\chref{recurrences} and \chref{master} --- the bound here is case 2 of the
master theorem. \chref{invariants} for the correctness argument.
\end{prereq}

\begin{keyterms}
divide and conquer $\bullet$ divide $\bullet$ conquer $\bullet$ combine
$\bullet$ merge $\bullet$ stability $\bullet$ in-place $\bullet$ auxiliary
space $\bullet$ cutoff $\bullet$ hybrid sort $\bullet$ top-down $\bullet$
bottom-up $\bullet$ branch misprediction
\end{keyterms}

\section{The Pattern}

\begin{definitionbox}[Divide and Conquer]
A \term{divide-and-conquer} algorithm has three steps:

\begin{tightnum}
  \item \textbf{Divide} the problem into subproblems that are smaller instances
        of the \emph{same} problem.
  \item \textbf{Conquer} the subproblems by solving them recursively; when they
        are small enough, solve them directly.
  \item \textbf{Combine} the subproblem solutions into a solution for the
        original.
\end{tightnum}

\smallskip
The running time is then $T(n) = a T(n/b) + D(n) + C(n)$, with $D$ and $C$ the
cost of dividing and combining --- which is exactly the form \chref{master}
solves.
\end{definitionbox}

\begin{conceptbox}[Where the work goes, and why it matters]
The three steps are not equally expensive, and which one carries the weight
distinguishes the two great sorting algorithms.

\smallskip
\textbf{Merge sort} divides trivially (compute the midpoint, $\bigO{1}$) and
combines expensively (merge two runs, $\bigTh{n}$). \emph{Easy split, hard
join.}

\smallskip
\textbf{Quicksort} (\chref{quicksort}) divides expensively (partition around a
pivot, $\bigTh{n}$) and combines trivially (the pieces are already in place,
$\bigO{1}$). \emph{Hard split, easy join.}

\smallskip
Both give $T(n) = 2T(n/2) + \bigTh{n}$ when the split is even. The difference
is that merge sort \emph{guarantees} the even split and quicksort only hopes
for it --- which is the whole of \chref{quicksort}.
\end{conceptbox}

\section{Merge Sort}

\dsfig{\aoaalg{\algname{Merge-Sort}$(A, \mathit{lo}, \mathit{hi})$}{
\If{$\mathit{hi} - \mathit{lo} < 2$}
  \State \Return
    \Comment{0 or 1 element: already sorted}
\EndIf
\State $\mathit{mid} \gets \lfloor (\mathit{lo}+\mathit{hi})/2 \rfloor$
\State \algname{Merge-Sort}$(A, \mathit{lo}, \mathit{mid})$
  \Comment{conquer the left}
\State \algname{Merge-Sort}$(A, \mathit{mid}, \mathit{hi})$
  \Comment{conquer the right}
\State \algname{Merge}$(A, \mathit{lo}, \mathit{mid}, \mathit{hi})$
  \Comment{combine}
}}{\algname{Merge-Sort}. The divide is one line of arithmetic; all the work is
in \algname{Merge}.}{mergesort}

\dsfig{\aoaalg{\algname{Merge}$(A, \mathit{lo}, \mathit{mid}, \mathit{hi})$}{
\State $i \gets \mathit{lo}$;\quad $j \gets \mathit{mid}$;\quad
       $k \gets \mathit{lo}$
\While{$i < \mathit{mid}$ \textbf{and} $j < \mathit{hi}$}
  \If{$A[i] \le A[j]$}
    \State $B[k] \gets A[i]$;\ \ $i \gets i+1$
      \Comment{$\le$ keeps it stable}
  \Else
    \State $B[k] \gets A[j]$;\ \ $j \gets j+1$
  \EndIf
  \State $k \gets k+1$
\EndWhile
\While{$i < \mathit{mid}$}
  \State $B[k] \gets A[i]$;\ \ $i \gets i+1$;\ \ $k \gets k+1$
\EndWhile
\While{$j < \mathit{hi}$}
  \State $B[k] \gets A[j]$;\ \ $j \gets j+1$;\ \ $k \gets k+1$
\EndWhile
\State copy $B[\mathit{lo} \mathinner{\ldotp\ldotp} \mathit{hi}]$ back into $A$
}}{\algname{Merge}: two sorted runs become one in a single pass. Line 3's
$\le$ rather than $<$ is what makes merge sort stable, and it is the only
difference.}{mergeproc}

\begin{theorem}[Merge is correct]\label{thm:merge}
If $A[\mathit{lo}\mathinner{\ldotp\ldotp}\mathit{mid}]$ and
$A[\mathit{mid}\mathinner{\ldotp\ldotp}\mathit{hi}]$ are each sorted, then
after \algname{Merge} the range
$A[\mathit{lo}\mathinner{\ldotp\ldotp}\mathit{hi}]$ is sorted and holds the
same multiset of values.
\end{theorem}

\begin{proof}
Invariant for the first loop: \emph{$B[\mathit{lo}\mathinner{\ldotp\ldotp}k]$
holds the $k - \mathit{lo}$ smallest elements of the two runs, in sorted order;
and $A[i]$ and $A[j]$ are the smallest remaining elements of their respective
runs.}

\smallskip
\textbf{Initialisation.} $k = \mathit{lo}$, so $B[\mathit{lo}
\mathinner{\ldotp\ldotp}\mathit{lo}]$ is empty and holds the 0 smallest
elements vacuously. $i$ and $j$ point at the fronts of two sorted runs, so they
are the smallest remaining in each.

\smallskip
\textbf{Maintenance.} The smallest remaining element overall is $\min(A[i],
A[j])$, since each run is sorted and $A[i], A[j]$ are their minima. The body
copies exactly that element into $B[k]$ and advances the corresponding index,
so $B$ gains the next-smallest element and the two pointers again indicate
their runs' minima.

\smallskip
\textbf{Termination.} The loop ends when one run is exhausted. The measure
$(\mathit{mid}-i) + (\mathit{hi}-j)$ strictly decreases each iteration and is
bounded below, so the loop is finite. The two tail loops append the remainder
of whichever run is left --- already sorted, and all of it $\ge$ everything in
$B$ by the invariant. Every element is copied exactly once and none is
overwritten before being copied, so the multiset is preserved.
\end{proof}

\begin{theorem}[Merge sort's running time]\label{thm:mergetime}
\algname{Merge-Sort} runs in $\bigTh{n \lg n}$ in the best, worst and average
cases, and uses $\bigTh{n}$ auxiliary space.
\end{theorem}

\begin{proof}
\algname{Merge} performs between $\min(n_{1},n_{2})$ and $n_{1}+n_{2}-1$
comparisons and exactly $n_{1}+n_{2}$ moves, so it is $\bigTh{n}$ regardless of
the data. The divide is $\bigO{1}$ and the two subproblems have sizes
$\lfloor n/2 \rfloor$ and $\lceil n/2 \rceil$. Hence
\[ T(n) = 2T(n/2) + \bigTh{n} \]
for \emph{every} input, with no case distinction. By \chref{master} with
$a = b = 2$ and $f(n) = n$, the watershed is $n^{\log_{2}2} = n = \bigTh{f(n)}$,
so case~2 applies and $T(n) = \bigTh{n \lg n}$. Since the recurrence holds for
every input, the bound is simultaneously best-, worst- and average-case. The
buffer $B$ holds at most $n$ elements, giving $\bigTh{n}$ auxiliary space.
\end{proof}

\begin{alertbox}[The bound is case-independent, the running time is not]
Theorem~\ref{thm:mergetime} says the \emph{order of growth} does not depend on
the input. It does not say the running time is the same for every input, and
Section~8.4 measures a factor of 3.9 between two inputs of the same size.

\smallskip
Keeping those two statements apart is the main intellectual work of this
chapter.
\end{alertbox}

\section{Measured: the Bound Holds, the Clock Does Not}

\begin{worked}[Three Input Shapes, $n = 2{,}000{,}000$]
\texttt{code/ch08\_merge.cpp} sorts the same two million elements arranged
three ways, counting comparisons exactly and timing the sort.

\rowcolors{2}{white}{lightbg}
\begin{center}
\begin{tabular}{@{}lrrrl@{}}
\toprule
\rowcolor{primary!15}
shape & time (ms) & comparisons & $/\,n\lg n$ & vs sorted \\
\midrule
sorted   &  73.1 & 20{,}769{,}984 & 0.496 & --- \\
reversed &  73.5 & 21{,}132{,}864 & 0.505 & $\times 1.006$ \\
random   & 285.7 & 39{,}348{,}350 & 0.940 & $\times 3.910$ \\
\bottomrule
\end{tabular}
\end{center}
($n \lg n = 41{,}863{,}137$.)

\smallskip
\textbf{First, the good news.} Compare this with \chref{cases}, where insertion
sort on these same three shapes varied by a factor of \textbf{6{,}600}. Merge
sort's comparison count varies by a factor of 1.9 and its time by 3.9. The
guarantee is real.

\smallskip
\textbf{Second, the comparison counts are exactly what theory says.} On sorted
input every merge exhausts its \emph{left} run first --- every element of the
left run is smaller than every element of the right --- so the comparison loop
performs exactly $n_{1}$ comparisons instead of about $n_{1}+n_{2}$. Half as
many. Measured: $0.496\, n\lg n$ against random input's $0.940\,n\lg n$, a
ratio of $1.89$.

\smallskip
\textbf{Third, the time ratio is 3.91 and the comparison ratio is 1.89.} Those
do not match. Roughly half the slowdown is not accounted for by work done, and
Section~8.5 is about finding out what it is.
\end{worked}

\begin{worked}[The Cutoff, Tuned]
Merge sort's recursion bottoms out at $n = 1$, which means a great many calls
doing almost nothing. \chref{role} measured insertion sort beating merge sort
below $n \approx 65$; so stop recursing at some cutoff and insertion-sort the
piece instead.

\rowcolors{2}{white}{lightbg}
\begin{center}
\begin{tabular}{@{}rrr@{}}
\toprule
\rowcolor{primary!15}
cutoff & time (ms) & vs none \\
\midrule
   0 & 275.5 & 1.000 \\
   4 & 257.1 & 0.933 \\
   8 & 243.1 & 0.882 \\
  16 & 232.5 & 0.844 \\
 \textbf{32} & \textbf{226.4} & \textbf{0.822} \\
  64 & 231.0 & 0.838 \\
 128 & 257.0 & 0.933 \\
 256 & 247.7 & 0.899 \\
\bottomrule
\end{tabular}
\end{center}

\smallskip
\textbf{The best cutoff on this machine is 32, and it is worth 18\%.} The curve
is shallow between 16 and 64 --- anything in that range gets most of the
benefit --- and turns back upward beyond it, as insertion sort's quadratic term
begins to bite.

\smallskip
\textbf{Note what this is and is not.} It is a constant-factor improvement: the
algorithm is $\bigTh{n\lg n}$ before and after, and \chref{asymptotics}'s
notation cannot see the difference. It is also the reason every production sort
in the world does it.

\smallskip
\textbf{And note that 32 is a fact about this laptop}, not about merge sort.
A different cache, a different element size or a different compiler moves it.
That is why the lab measures it rather than quoting it.
\end{worked}

\begin{worked}[Four Implementations of One Bound]
\rowcolors{2}{white}{lightbg}
\begin{center}
\begin{tabular}{@{}lrr@{}}
\toprule
\rowcolor{primary!15}
implementation & time (ms) & relative \\
\midrule
top-down, no cutoff & 279.2 & 1.00 \\
top-down, cutoff 32 & 225.1 & 0.81 \\
bottom-up           & 263.3 & 0.94 \\
\texttt{std::sort}  & 169.1 & 0.61 \\
\bottomrule
\end{tabular}
\end{center}

\smallskip
All four are $\bigTh{n\lg n}$ and they span a factor of 1.65. Bottom-up merge
sort --- same recurrence, no call stack --- is 6\% faster than top-down, which
is less than most people expect; the recursion was never the expensive part.

\smallskip
\texttt{std::sort} wins by 39\% because it is not merge sort at all. It is
introsort: quicksort (\chref{quicksort}) with an insertion-sort cutoff and a
heapsort fallback, sorting in place and so touching half as much memory. The
comparison is here to make a point that $\bigTh{\cdot}$ cannot: \textbf{within
one order of growth there is still a factor of nearly two to be had, and it is
won by caring about memory.}
\end{worked}

\section{Where the Missing Factor of Two Went}

The gap in Section~8.4 was 3.91 in time and 1.89 in comparisons. Something
costs time without being a comparison.

\begin{conceptbox}[A hypothesis, and how to test it]
\textbf{Hypothesis.} The culprit is the branch \texttt{if (A[i] <= A[j])} in
the merge loop. On sorted input it goes the same way for an entire run and the
processor's branch predictor gets it right every time. On random input it is a
coin flip, mispredicted about half the time, at a cost of roughly fifteen
cycles each.

\smallskip
\textbf{This is exactly the third lie of \chref{role}'s RAM model} --- the one
that says instructions execute one at a time --- coming back to charge for
itself.

\smallskip
\textbf{The test.} Rewrite the merge loop so that the choice is arithmetic
rather than a branch:
\begin{center}
\texttt{bool left = A[i] <= A[j];} \\
\texttt{B[k++] = left ? A[i] : A[j];\quad i += left;\quad j += !left;}
\end{center}
This compiles to a conditional move. \emph{Same comparisons, same output, no
branch to mispredict.} If the hypothesis is right, the gap should close. If it
is wrong, nothing will change and we will have learnt that too.
\end{conceptbox}

\begin{worked}[The Test, Run]
\rowcolors{2}{white}{lightbg}
\begin{center}
\begin{tabular}{@{}lrr@{}}
\toprule
\rowcolor{primary!15}
merge variant & time (ms) & vs branchy random \\
\midrule
branchy, random input    & 272.9 & 1.000 \\
branchless, random input & 152.0 & \textbf{0.557} \\
branchy, sorted input    &  72.1 & 0.264 \\
\bottomrule
\end{tabular}
\end{center}

\smallskip
\textbf{Removing the branch cut the time on random input by 44\%.} The
hypothesis is confirmed.

\smallskip
\textbf{And now the gap decomposes cleanly:}
\[ \underbrace{3.78}_{\text{random/sorted, branchy}}
   \;=\; \underbrace{1.89}_{\text{more comparisons}}
   \;\times\; \underbrace{2.00}_{\text{branch misprediction}} \]
Measured: random/sorted was $3.78\times$ with the branch and $2.11\times$
without it. So $3.78 / 2.11 = 1.79$ is attributable to the branch, and the
residual $2.11$ matches the comparison-count ratio of $1.89$ to within 12\%.

\smallskip
\textbf{Two lessons.} First, a factor of two was hiding in a line of code that
$\bigTh{\cdot}$ cannot see and that no amount of algorithm analysis would have
revealed --- only measurement found it, and only a designed experiment
explained it. Second, the branchless version is not simply better: it does the
\emph{same} comparisons, so on sorted input it is slower than the branchy one,
which gets its branches free. Which version wins depends on the data, and that
is a sentence this book has now written in four different chapters.
\end{worked}

\begin{pitfall}
\begin{tight}
  \item \textbf{Reading ``case-independent bound'' as ``case-independent
        time''.} Measured, the same sort on the same $n$ varied by 3.9.
  \item \textbf{Allocating the buffer inside the recursion.} One buffer for the
        whole sort, allocated once. \chref{asymptotics}'s lab nearly reported
        nonsense for exactly this reason.
  \item \textbf{Using $<$ instead of $\le$ in the merge.} It still sorts; it is
        no longer stable, and stability is often why merge sort was chosen.
  \item \textbf{Claiming merge sort is in-place.} It needs $\bigTh{n}$
        auxiliary space. In-place merging exists and is slower and much harder.
  \item \textbf{Quoting a cutoff from a book.} The best value here was 32, on
        this machine, for 4-byte elements. Measure your own.
  \item \textbf{Tuning the cutoff before fixing the algorithm.} 18\% from a
        cutoff is worth having; it is not worth having instead of a better
        order of growth.
  \item \textbf{Assuming recursion is the overhead.} Bottom-up was only 6\%
        faster. The memory traffic is the cost, not the call stack.
\end{tight}
\end{pitfall}

\begin{lab}[Merge Sort, Tuned and Explained]
Reproduces all four tables. Expect thirty minutes; it sorts two million
elements many times.

\begin{lstlisting}[language=C++]
// ch08_merge.cpp -- the divide-and-conquer pattern, and what tuning buys.
// Build: cl /O2 /EHsc /std:c++17 ch08_merge.cpp     (see Appendix A)

// --- 1. The merge, which is the whole idea -------------------------
// Two sorted runs become one in a single pass: look at both fronts,
// take the smaller. Linear in the total length, and the only place
// the algorithm compares anything.
static void merge(std::vector<int>& a, std::vector<int>& buf,
                  size_t lo, size_t mid, size_t hi, bool count) {
    size_t i = lo, j = mid, k = lo;
    while (i < mid && j < hi) {
        if (count) ++comparisons;
        buf[k++] = (a[i] <= a[j]) ? a[i++] : a[j++];
    }
    while (i < mid) buf[k++] = a[i++];
    while (j < hi)  buf[k++] = a[j++];
    std::copy(buf.begin() + lo, buf.begin() + hi, a.begin() + lo);
}

// --- 2. Top-down: the recurrence written out as code ---------------
// T(n) = 2T(n/2) + Theta(n), which Chapter 6 settles as Theta(n lg n).
// The cutoff is the one tuning knob: below it, insertion sort's
// smaller constant wins, exactly as Chapter 1 measured.
static void msort(std::vector<int>& a, std::vector<int>& buf,
                  size_t lo, size_t hi, size_t cutoff, bool count) {
    if (hi - lo < 2) return;
    if (hi - lo <= cutoff) { insertion(a, lo, hi); return; }
    size_t mid = lo + (hi - lo) / 2;
    msort(a, buf, lo, mid, cutoff, count);
    msort(a, buf, mid, hi, cutoff, count);
    merge(a, buf, lo, mid, hi, count);
}

// --- 3. The same merge with the branch removed ---------------------
// Same comparisons and same output; the choice becomes arithmetic
// that compiles to a conditional move, so there is no branch to
// mispredict. Used only by part 7, to test what the gap is made of.
static void merge_bl(std::vector<int>& a, std::vector<int>& buf,
                     size_t lo, size_t mid, size_t hi) {
    size_t i = lo, j = mid, k = lo;
    while (i < mid && j < hi) {
        bool take_left = a[i] <= a[j];
        buf[k++] = take_left ? a[i] : a[j];
        i += take_left ? 1 : 0;
        j += take_left ? 0 : 1;
    }
    while (i < mid) buf[k++] = a[i++];
    while (j < hi)  buf[k++] = a[j++];
    std::copy(buf.begin() + lo, buf.begin() + hi, a.begin() + lo);
}

// --- 4. The bound does not care what the input looks like ----------
// Chapter 2 found insertion sort varying by 6,600x across these same
// three shapes. Merge sort should vary by almost nothing.

// --- 5. What is the cutoff worth, and where is it? -----------------
// The best value is a property of this machine, not of the algorithm,
// which is exactly why it has to be measured rather than looked up.

// --- 6. Top-down against bottom-up, and against the library --------
// Same recurrence, same Theta. Everything that differs is constant
// factors -- recursion overhead, memory order, and in std::sort's
// case a different algorithm entirely (Chapter 9).

// --- 7. Why is random input 4x slower than sorted input? -----------
// Part 4 found 3.9x in time but only 1.9x in comparisons, so half the
// slowdown is NOT work done. Hypothesis: the branch in the merge
// loop, perfectly predicted on sorted input and a coin flip on
// random. Replace it with arithmetic and see whether the gap closes.
\end{lstlisting}

\textbf{What to notice.} Part 7 is the shape of a real investigation: a
measurement that does not fit, a hypothesis that names a mechanism, and an
experiment whose two possible outcomes would tell you different things. The
answer --- 0.557 --- was not knowable in advance.

\smallskip
\textbf{Then try to break it.} Run the branchless merge on \emph{sorted} input
and predict the result first. It should be \emph{slower} than the branchy one,
because a perfectly predicted branch costs almost nothing while a conditional
move always costs something. If your prediction was ``faster'', the thing to
revise is the belief that branchless code is simply better.
\end{lab}

\begin{checkpoint}
\begin{tightnum}
  \item Why is merge sort's recurrence $T(n) = 2T(n/2) + \bigTh{n}$ valid for
        every input, while insertion sort needs three separate analyses?
  \item \algname{Merge} performs between $\min(n_{1},n_{2})$ and
        $n_{1}+n_{2}-1$ comparisons. Describe an input achieving each bound,
        and relate your answer to the measured 0.496 and 0.940.
  \item The measured best cutoff was 32 and \chref{role}'s measured crossover
        was 65. Give a reason the two numbers differ.
\end{tightnum}
\end{checkpoint}

\begin{chaptersummary}
\begin{tight}
  \item Divide and conquer: divide into smaller instances of the same problem,
        conquer recursively, combine. The cost is $aT(n/b) + D(n) + C(n)$.
  \item Merge sort has an easy split and a hard join; quicksort the reverse.
        Merge sort \emph{guarantees} the even split, which is why its bound
        holds in every case.
  \item $T(n) = 2T(n/2) + \bigTh{n}$ for every input, so by case 2 of the
        master theorem merge sort is $\bigTh{n\lg n}$ best, worst and average,
        with $\bigTh{n}$ auxiliary space.
  \item Measured at $n = 2{,}000{,}000$: comparisons were $0.496\,n\lg n$ on
        sorted input and $0.940\,n\lg n$ on random, and times were 73 ms and
        286 ms. The bound is case-independent; the clock is not.
  \item An insertion-sort cutoff of 32 was worth 18\% on this machine. The best
        value is a property of the machine and must be measured.
  \item The $3.78\times$ sorted-to-random time gap decomposes as $1.89\times$
        (more comparisons) $\times$ $2.0\times$ (branch misprediction).
        Replacing the branch with a conditional move cut random-input time by
        44\%, which confirmed the mechanism.
  \item \texttt{std::sort} beat all four merge sorts by 39\% --- the same order
        of growth, won back by sorting in place.
\end{tight}
\end{chaptersummary}

\begin{reviewq}
  \item \textbf{(MCQ)} Merge sort's auxiliary space requirement is:
        (a)~$\bigTh{1}$; (b)~$\bigTh{\lg n}$; (c)~$\bigTh{n}$;
        (d)~$\bigTh{n\lg n}$.
  \item \textbf{(MCQ)} Merge sort's worst-case running time is:
        (a)~$\bigTh{n}$; (b)~$\bigTh{n\lg n}$; (c)~$\bigTh{n^{2}}$;
        (d)~dependent on the input order.
  \item \textbf{(Short)} State the three steps of divide and conquer and apply
        them to merge sort, naming the cost of each \emph{(CLO-6)}.
  \item \textbf{(Short)} Explain why merge sort is stable, and which single
        character in \algname{Merge} is responsible.
  \item \textbf{(Short)} A colleague reports that their merge sort is twice as
        slow on one dataset as on another of the same size, and concludes the
        implementation is buggy. Give two innocent explanations.
  \item \textbf{(Applied)} Derive the number of comparisons \algname{Merge}
        performs on already-sorted input, and use it to predict the
        $0.496\,n\lg n$ measured above.
  \item \textbf{(Applied)} Design an experiment to find the best insertion-sort
        cutoff for sorting 64-byte records rather than 4-byte integers. Predict
        whether the best cutoff will be larger or smaller than 32, and say why.
\end{reviewq}
